\paragraph{The retirement contract.}
\label{sec:retirement}
Fix a session $(\beta,r)$ and write $\Delta=\Delta_{i_{\beta,r}}$,
using its immutable voting snapshot. Closure is a Session Manager
obligation independent of whether a program exists or has entered.
The manager persistently stores authenticated $\tagg{Closed}_v[\beta]$
and its complete source evidence, not just a hash or retrieval promise.
The regional control owner keeps those bytes available for the receive
handler below; a pending storage copy does not suspend that owner.
This branch-wide record immediately
prohibits every old-$\beta$ initial entry and fresh core action.
It is separate from a program's six phases and from the later
per-session complete-close event defined below.

For each known fixed session snapshot in which $v$ is a member, the
manager retains a control record even in Waiting, Ready, Completed or
Terminated, or when no worker exists. The record includes the historical
fixed boundary, source region, immutable authority and membership,
\texttt{close\_requested}, any first complete-close time $c_v$, and
pending relays. Active/Frozen programs keep their core phase while
draining; only genuine Quit or terminal Completion changes that phase.
Force is neither a quit proof nor a terminal receipt.

\paragraph{Complete-receipt dissemination.}
\label{sec:retirement-receipt-relay}
The following obligation applies during normal running as well as
retirement. The first time a correct regional member or its local
manager/collector holds a complete valid receipt of its own session,
it multicasts the receipt and all validation dependencies to the
session's fixed membership. In particular, a terminal receipt is
multicast no later than its local Completion event. Deduplication by
session and DATA hash suppresses repeated encodings or signer sets,
not this first complete relay. The relay and any pending delivery are
preserved across Completed and Terminated by the persistent owner.
Recipients process an own-session terminal receipt under Completion
even while Frozen or closing, independently of its generation and
before a simultaneous timer. A program that has not entered caches
it; a Terminated program archives it without restarting. This rule
forwards existing signatures and introduces neither a new vote nor a
certificate. The intra-region complete-message bound applies to the
actual relay; it is not a bound on first forming a receipt.

\paragraph{Manager closure and \texttt{RequestClose}.}
The following handler runs at the manager or persistent control owner,
not only at an entered engine.
\begin{enumerate}[label=(\roman*),itemsep=2pt]
\item \emph{Record public closure.} Verify complete evidence that the
canonical prelude of an authenticated finalized Espresso block closed
$\beta$. Retain its triggering header, applicable complete control and
namespace history from a validated checkpoint, and closure and
retained-prefix context. Check the existing MSIG/AUTO rules; a pending
MSIG, OFF, unfinalized header or bare timestamp is insufficient.
Atomically persist $\tagg{Closed}_v[\beta]$, prohibit fresh actions for
that version and notify all its local control records, regardless of
phase. This consumes an established closure; it does not execute Force
or select an anchor a second time.
\item \emph{Disseminate for every known snapshot.} For each known fixed
session $(\beta,r)$ of this region, combine the retained closure proof
with its immutable authority, membership and historical admission
context. This is a \defterm{complete close bundle}: another member can
validate both closure and the exact recipient snapshot without a prior
program, live-entry permission or an unprovided cache. On first complete
processing, perform RequestClose atomically: set
\texttt{close\_requested}, record $c_v$ as the current time, and multicast
the bundle to the fixed membership. This includes Waiting, Ready,
Active, Frozen, Completed and Terminated records. A recipient validates
the bundle, persists Closed and performs the same idempotent operation;
no worker or engine entry is created. Repeated bundles retain $c_v$.
\item \emph{Retain unresolved dissemination.} If the fixed snapshot or
its dependencies are not yet known, retain the authenticated Closed
record and a pending relay obligation. The entry gate is already shut,
but $c_v$ is not assigned and no $4\Delta$ clock is invented.
Whenever any old-pair Ready, proposal, authority, history or other packet
provides the missing authenticated snapshot context, the manager or
archive completes and multicasts the close bundle in that receive
handler, and replies with it to the packet's sender. Do this before
ordinary BFT processing, rejecting a live-entry request, discarding a
duplicate, or returning a Terminated/tombstone result. Historical
context is validated without requiring that the old version remain
open. Complete proposals from correct voters carry or resolve this
context within the existing complete-message bound. This handler is
mandatory even if the old worker was never created, its ReadyBundle
is refused, or a delayed callback arrives after exit or restart.
\item \emph{Gate new work.} Closed prohibits every old-version
DATA/DRAIN proposal, NEWVIEW, DATA/DRAIN/ANCHOR vote, new commit timer,
progress check producing BLAME, term entry and generation entry.
The checks cover initial Enter, delayed authority/recovery callbacks,
terminal-QC rotation and buffered-message processing. Neither a new
STATUS nor recreation of a local worker bypasses the persistent gate.
Already dispatched messages are not withdrawn. New-session preparation
and entry use their own pair and need not wait for retirement.
\item \emph{Drain existing work.} Preserve existing commit timers and
their original deadlines; observing Closed alone does not cancel a
timer or release its evidence-processing owner. Continue complete
validation, comparison with retained signed material, and evidence-first
processing. A genuine current-generation quit proof still runs Quit,
broadcasts its proof, cancels its timers and preserves commitments;
it cannot reopen the gate. A surviving timer commits its block and
ancestors. For each DATA thereby committed, its Protocol~1 FINAL goes
both to the fixed regional membership and to the publication collector,
preserving \cref{def:receipt}. A partial FINAL does not invoke the
complete-result guarantee of \cref{ass:publish}.
\item \emph{Keep evidence duties.} Continue collecting old votes, QCs,
FINAL messages and receipts. Multicast complete own-session receipts
as specified above, and offer every complete result to the Publish
Layer. A valid terminal receipt still runs genuine Completion,
possibly changing the phase to Completed and canceling timers; an
existing drain deadline remains in force. Neither receipt arrival nor
later ON reopens the branch. Keep data, control-relay and publication
responsibilities until their durable transfer.
\end{enumerate}
All closure relays use existing authenticated source evidence and
signatures, not new BFT votes or certificates. The manager/control owner
continues transporting them after worker termination. A known snapshot
is never forgotten merely because its program exits. Recovery from a
crash restores Closed, session close times, pending relays and owners
before dispatching callbacks or signatures. Messages, including complete
closure, quit and terminal-receipt evidence, are processed before a
simultaneous timer, entry or retirement event. No receipt, lock, vote,
committed block or provenance is erased by this handler.

\paragraph{\texttt{FinishClose}.}
For a program Active or Frozen when it first learns closure, first
require its complete per-session RequestClose event. Set
\[
 d_v=c_v+4\Delta.
\]
An earlier Closed flag without a complete-close event is not a substitute
for $c_v$ and never authorizes early release of timer/evidence processing.
FinishClose may mark it Terminated only at a time $t\ge d_v$ when:
\begin{enumerate}[label=(\roman*),itemsep=2pt]
\item every existing commit timer has executed or was canceled by
genuine Quit or Completion, and no queued action can pass a fresh-action
gate;
\item persistent storage owns the exit, Closed and per-session close records,
committed log, required signed evidence and validation dependencies,
receipts, anchor data and recoverable state, without releasing the last
required copy;
\item a persistent owner has accepted pending closure-context discovery,
complete-close relays and replies, complete-receipt multicasts,
complete-result offers, publication retries and data-service duties, and each locally complete result has been offered to the
Publish Layer; and
\item the dispatcher routes later messages for $(\beta,r)$ to those
owners. They must run the manager-level closure handler even for an
unrecognized old pair, and may collect or relay existing signatures,
validate late receipts, serve data and continue publication, but cannot restart the
engine, execute its BFT timers or issue fresh BFT signatures for it.
\end{enumerate}
The ownership transfer may be local to the same correct replica; it
requires no additional agreement or certificate. The receiving control
owner must already hold the required evidence and receive route before
the old owner releases them; an unresolved retrieval task is not such
a handoff. Only after these
guards hold may the production/recovery worker and its timer machinery
be released. A late terminal receipt records archival evidence and does
not change a Terminated program back to Completed or Active.

A definitively invalid Waiting or Ready program has no regional timers
to drain. It may release its worker without waiting $4\Delta$ only after
its durable exit record and all manager/control duties above have an
owner. In particular, learning Force before its first ReadyBundle never
permits silent cancellation: known snapshots trigger immediate
RequestClose; unknown snapshots leave pending relays that the refused
ReadyBundle or a later proposal must resolve. Already-Completed programs
have performed the original timer cancellation and may release workers
after the same responsibility transfer. They likewise retain all
control-relay duties. A temporary candidate-eligibility change alone
need not terminate preparation. Termination never deletes a commitment,
removes DATA from $G$, or authorizes stake withdrawal.

\begin{proposition}[No silent closure in a vote window]
\label{prop:retirement-observable-close}
Let a correct member vote for and completely forward a valid proposal
of $(\beta,r)$ at time $s$. If a correct member's matching vote or entry
is suppressed by its authenticated Closed record during processing of
that proposal by $s+\Delta$, a correct member processes and multicasts
a complete close bundle for this session by $s+\Delta$. Thus, writing
$c=\min_{v\ {\rm correct}}c_v$ over all phase-independent complete-close
events (with $c=+\infty$ if none occurs), such suppression implies
$c\le s+\Delta$.
\end{proposition}
\begin{proof}
The forwarded proposal and its historical authority/admission dependencies
reach the member's manager or archive within $\Delta$, even if it has
no worker or an exit tombstone. These dependencies identify the fixed
regional snapshot. If Closed was already recorded, the receive handler
combines them with its retained closure evidence and sends the complete
bundle before refusing ordinary processing. If closure is learned after
that packet but by the time it suppresses a vote, the snapshot is already
known and closure's own handler performs the multicast immediately.
Either way some $c_v\le s+\Delta$. No new initial entry, executable-state
restoration or remote membership discovery is required for this control
relay. Restart and duplicate paths use the same persistent handler.
\end{proof}

\begin{proposition}[Closure propagation and honest-message drain]
\label{prop:retirement-drain}
Under \cref{ass:network}, let $c=\min_{v\ {\rm correct}}c_v$ be the first
complete per-session RequestClose time among this snapshot's correct
members in any phase, including its manager/archive owners. Then every correct member has closed by $c+\Delta$.
Every surviving honest commit timer expires by $c+3\Delta$, and the
regional FINAL messages generated by those timer events reach every
correct member by $c+4\Delta$. Consequently $d_v=c_v+4\Delta$ is a
conservative local drain threshold, not a deadline for finishing all
retirement duties.
\end{proposition}
\begin{proof}
The first correct relay includes closure and the fixed snapshot's full
validation context. Its delivery and processing are covered by the
existing regional bound even when no receiving worker exists. Thus
$c\le c_v\le c+\Delta$ and every old-version fresh-action gate is shut
by $c+\Delta$. A member may have set Closed still earlier; that can only
stop its votes earlier. A correct DATA/DRAIN vote at time
$s=\tau_v(P)$ starts its only timer, with expiry $s+2\Delta$.
The close gate gives $s\le c_v\le c+\Delta$; no callback starts a
replacement timer. Thus its expiry is at most $c+3\Delta$.
Quit and Completion can cancel a timer but cannot postpone it.
The FINAL messages actually generated at surviving expiries are
delivered regionally within another $\Delta$. Finally $c_v\ge c$
implies $d_v\ge c+4\Delta$. Message-first processing covers equality.
\end{proof}

The bound starts at complete per-session closure processing and its
actual regional relay. It does not start at an earlier Closed observation
whose session snapshot was unknown, a control signature, the triggering Espresso timestamp, or the
start of unbounded historical catch-up. It does not add a new source
visibility assumption. Three local $\Delta$ intervals cover the last
honest timer event but not its FINAL delivery: a member can vote just
before $c+\Delta$, commit just before $c+3\Delta$, and deliver FINAL
only just before $c+4\Delta$.

Four $\Delta$ intervals do not bound subsequent receipt relays,
Byzantine message release, a storage copy, $T_{\rm pub}$, $T_E$, or
successful EXPORT acceptance. Their outstanding duties pass to the
persistent owners; no bounded storage size or transfer latency is
claimed. Elapsed grace alone therefore never authorizes FinishClose.

\begin{proposition}[Conditional regional safety during retirement]
\label{prop:retirement-safety}
Assume the stated regional model, the baseline regional validation
and DATA/DRAIN soundness premises, and the complete-message Ready/Enter
installation contract of \cref{sec:ready}. With the phase-independent manager closure handler,
complete-receipt multicast and retirement rules above, correct regional committed logs
remain prefix-comparable and irrevocable. The external prefix argument
of \cref{sec:safety-argument} is preserved. This does not assert a
protecting QC for every late surviving timer, old-session progress,
or complete-receipt formation.
\end{proposition}
\begin{proof}
We verify the prerequisites of the timer/lock argument explicitly;
baseline safety is not used to infer an unproved universal-QC property.
Fix a correct direct timer for DATA or DRAIN block $B$ in generation
$g$, voted at $s$, that actually expires at $s+2\Delta$ without
cancellation. All comparisons below concern this same session.
Let $e_0$ be its first correct initial-entry time. The specified Ready
contract enters every nonclosed correct member by $e_0+\Delta$.
An own-session terminal receipt requires a correct FINAL after a timer,
whose earliest possible expiry is $e_0+2\Delta$. Hence such a receipt
cannot evade Completion by arriving only at an unentered cache of a
correct member that will run the chosen timer.

\emph{Quit and Completion cannot suppress its early vote window.}
If any correct member first quits $g$ at $a\le s+\Delta$, its complete
quit proof reaches the committer by $a+\Delta\le s+2\Delta$ and
cancels the timer. Thus any first correct Quit satisfies
$a>s+\Delta$. If any correct member Completes at $b\le s+\Delta$,
its complete terminal receipt is multicast by $b$ and reaches the
committer by $b+\Delta\le s+2\Delta$. The committer then executes
Completion: it commits the terminal DATA and ancestors and cancels
this timer, including a DRAIN timer. Again this contradicts the chosen
surviving direct timer. No correct Completion can therefore occur at
or before $s+\Delta$. This is an operational cancellation argument;
it does not assume that receipt-triggered commitments are already
known to be prefix-comparable.

\emph{The proposal is processable within its regional window.}
For $g=0$, the self-contained Ready relay and installation contract
in \cref{sec:ready} supplies entry by $s+\Delta$; this assertion is
conditional on that precise contract, not a bound on state restoration
before any member is Ready. For $g>0$, let $a^-$ be the first correct
Quit of $g-1$, and let $e$ be the first correct entry into $g$.
The original recovery rule gives $e\ge a^-+\Delta$ and $s\ge e$.
The complete quit relay reaches every correct member by $a^-+\Delta$;
its own mandatory wait has therefore ended by
$a^-+2\Delta\le e+\Delta\le s+\Delta$. The forwarded complete
proposal supplies the finalized authority and remaining admission
history. Hence a receiver with no Quit, Completion or close preventing
entry can process it in $g$ by $s+\Delta$. Applying the same
complete-proof argument inductively to earlier generations supplies
any predecessor entry needed to process the quit proof. This uses
the existing recovery waits and complete validation bundles, not a
new simultaneous authority-visibility time.

\emph{Early and late windows.}
Let $c$ be the first complete per-session RequestClose event among
correct manager/control owners in any phase. A prior Closed flag or
unentered/Terminated record is not excluded merely by definition of $c$.
Instead, \cref{prop:retirement-observable-close} shows that if it
suppressed a matching vote or entry by $s+\Delta$, then $c\le s+\Delta$.
Consequently, if $s+\Delta<c$, no closure suppresses that vote or entry.
The preceding cancellation arguments also exclude Quit and Completion
through that window. Earlier-term processing and prior indirect
commitment do not suppress that vote either. An incompatible retained
history would expose a genuine quit proof in this window and contradict
timer survival. Every correct member therefore has cast, or now casts,
the matching DATA/DRAIN vote by $s+\Delta$. Their total weight is at least
$W-F>F$, so the votes form $B$'s protecting QC by $s+2\Delta$ at
all members that can still take a later recovery entry. A member that
Completed after the vote window cannot take such an entry; it need
not be assigned an unused recovery lock.

Any correct entry into $g+1$ must obey
\[
 e_v(\beta,r,g+1,0)\ge a_v+\Delta
 \ge a+\Delta>s+2\Delta.
\]
In the early case, every such entrant therefore knows the protecting
QC before its lock snapshot. The highest-QC rule gives the baseline
recovery-lock invariant. Closure preserves collecting those already
sent votes and QCs, and may only remove a later entry. In the late
case $s+\Delta\ge c$, the timer expires at or after $c+\Delta$.
The complete snapshot-bearing relay has closed every correct member's
entry gate by then, including previously unentered and archived members; the strict displayed
inequality prohibits every later-generation entry and thus its
correct signatures. A later-generation QC needs a correct signer,
so no recovery can exploit a missing protecting QC in this case.

\emph{Direct timers and receipt-based commits.}
In one generation, complete signed histories either extend one
another or expose the original admissible equivocation proof.
A correct voter forwards its complete history. A correct signer of
an incompatible history either processes $B$ first and cannot vote,
or forwards its history before processing $B$, no later than
$s+\Delta$. That history reaches the committer by expiry and exposes
the proof. Genuine Quit, or Completion before expiry, would cancel
the timer. Thus no conflicting same-generation correct direct timer
can survive; the same argument excludes a conflicting QC, including
one assembled later from old votes.

Across generations, a surviving lower-generation direct timer cannot
be in the late case if a correct higher-generation vote exists.
Its early-case QC is known before every later live lock snapshot.
The baseline ranked-QC descent now has its actual prerequisites:
correct ANCHOR voters meet those locks; DATA/DRAIN QCs inherit a
lower-ranked parent justification. It excludes a later QC that does
not extend the committed block, and the parent of a later direct
proposal consequently extends it. These are precisely the
protecting-QC, recovery-lock and unique-extension steps of Part~I,
with the explicit vote-window checks above and the baseline
DATA/DRAIN validation rules. Direct commitments are comparable, and
ancestor commitments inherit that property.

Finally a valid receipt contains a correct FINAL signer. That signer
sent FINAL only after a timer committed its named DATA, directly or
as an ancestor; committing from another receipt emits no FINAL.
Consequently every receipt-based commit, including Completion, traces
to a causally earlier direct timer commitment of that DATA or a
descendant. It adds an already protected prefix. This provenance
argument does not take the safety of receipt-based commits as its own
premise. A timer canceled by Completion contributes no additional
direct commitment, particularly no commitment of its unexpired DRAIN.
No retirement step removes a prior commitment.

The public interpreter still preserves all of $G$ at its canonical
Force boundary and rejects further old-version DATA. Retirement and
receipt relay neither select a retained anchor nor authorize an
append, so those external transitions are unchanged.
\end{proof}

The universal protecting-QC lemma from Part~I cannot simply be copied:
its proof requires matching votes from all correct members, whereas a
closed member now retains a proposal without voting. For example, with
three unit-weight members and $F=1$, let correct $H_0$ close at $c=0$
and correct $H_1$ receive closure at $\Delta$. A Byzantine leader gives
a valid proposal to $H_1$ at $0<s<\Delta$. $H_1$ votes, forwards and
later commits; $H_0$ retains it without voting. If the Byzantine member
withholds its vote and FINAL, only $H_1$'s FINAL exists. Its weight is
not greater than $F$: neither a proposal QC nor a complete receipt need
ever form. Retirement preserves this local commitment and evidence
without calling it Completion. The Publish Layer's complete-result
premise does not repair the missing signature.

\paragraph{Force before the first ReadyBundle.}
Suppose $H_0$ enters at $e_0$ and relays its complete ReadyBundle, while
$H_1$ learns Force before that bundle arrives. $H_1$ immediately stores
Closed. If its snapshot is already known it multicasts closure at once;
otherwise the ReadyBundle supplies the historical snapshot by
$e_0+\Delta$ and triggers the same multicast before its entry refusal.
The complete close evidence therefore reaches $H_0$ by $e_0+2\Delta$,
regardless of $H_0$'s later direct Espresso retrieval. Control processing
wins equality with its initial leader wait. If a proposal is the first
packet to reveal the snapshot instead, \cref{prop:retirement-observable-close}
gives the required $c\le s+\Delta$ and the late-window argument blocks
later recovery. Worker termination, a stale callback and restart cannot
remove this relay path. This rules out a silently absent member being
used to classify an unprotected timer as an early-window commitment.
